\documentclass[12pt]{article}
\usepackage[utf8]{inputenc}
\usepackage[T1]{fontenc}
\usepackage{hyperref}
\title{Temuan Language Resources for K-12 Curriculum Development}
\author{Compiled by Claude Sonnet 5 (claude-sonnet-5) for Andr\'as Kornai}
\date{}
\begin{document}
\maketitle

\section{Introduction}
Temuan (ISO 639-3 tmw; Glottocode temu1239) is an Austronesian language of the Orang Asli in Peninsular Malaysia, with about 23{,}000 speakers. Schooling is in Malay. It is mostly written in the Latin alphabet, and no standard orthography exists.

\section{Primary Educational Resources}
No Temuan primers or textbooks were found. One study of Orang Asli language education covers Temuan children:
\begin{itemize}
\item Language Education for Orang Asli Children in Malaysia (University of Malaya): \url{https://eprints.um.edu.my/35783}
\end{itemize}

\section{Linguistic and Reference Materials}
\begin{itemize}
\item Wikipedia, Temuan language. It mentions a Temuan-Malay-English lexicon of about 3{,}200 entries by Looi, Baer and Mohamad (2018, Center for Orang Asli Concerns), using a non-standard phonetic transcription: \url{https://en.wikipedia.org/wiki/Temuan_language}
\item Wikipedia, Temuan people: \url{https://en.wikipedia.org/wiki/Temuan_people}
\end{itemize}

\section{Religious and Community Resources}
No Temuan Bible translation was found in search results.

\section{Digital and Audio Resources}
No audio or software resources were found in search results.

\section{Curriculum Development Context}
Literacy among Orang Asli children is reported to be well below the national average, and teaching is in Malay. A mother-tongue curriculum would need an agreed orthography. The 2018 lexicon could serve as a starting vocabulary once it is checked with speakers.
\end{document}